\section*{Exercises on \S\arabic{exercisegroup}}
\phantomsection
\addcontentsline{toc}{section}{Exercises on \protect\S\arabic{exercisegroup}}

\begin{enumerate}
\def\labelenumi{\arabic{enumi})}
\tightlist
\item
  \begin{enumerate}
  \def\labelenumii{\alph{enumii})}
  \tightlist
  \item
    Let H be a set of convex functions on a compact interval \([a, b] \subset R\) ; suppose that the sets \(\mathrm{H}(a)\) and \(\mathrm{H}(b)\) are bounded above in R and that there exists a point c such that a \textless{} c \textless{} b and that \(\mathrm{H}(c)\) is bounded below in R; show that H is an equicontinuous set on {]}a, b{[} (Gen.~Top., X, p.~283).
  \end{enumerate}
\end{enumerate}

\begin{enumerate}
\def\labelenumi{\alph{enumi})}
\setcounter{enumi}{1}
\tightlist
\item
  Let H be a set of convex functions on an interval \(I \subset R\) , and let F be a filter on H which converges pointwise on I to a function \(f_{0}\) ; show that F converges uniformly to \(f_{0}\) on every compact interval contained in I.
\end{enumerate}

\begin{enumerate}
\def\labelenumi{\arabic{enumi})}
\setcounter{enumi}{1}
\item
  Show that every convex function \(f\) on a compact interval \(\mathrm{I} \subset \mathbf{R}\) is the limit of a decreasing uniformly convergent sequence of convex functions on \(\mathrm{I}\) which admit a second derivative on \(\mathrm{I}\) (first consider the function \((x - a)^+\) , and approximate \(f\) by the sum of an affine linear function and a linear combination \(\sum c_j(x - a_j)^+\) with coefficients \(c_j \geqslant 0\) ).
\item
  Let \(f\) be a convex function on an interval \(I \subset \mathbf{R}\) .
\end{enumerate}

\begin{enumerate}
\def\labelenumi{\alph{enumi})}
\item
  Show that if \(f\) is not constant it cannot attain its least upper bound at an interior point of \(I\) .
\item
  Show that if I is relatively compact in R then f is bounded below on I.
\item
  Show that if \(\mathbf{I} = \mathbf{R}\) and \(f\) is not constant, then \(f\) is not bounded above on \(\mathbf{I}\) .
\end{enumerate}

\begin{enumerate}
\def\labelenumi{\arabic{enumi})}
\setcounter{enumi}{3}
\item
  For a function \(f\) to be convex on a compact interval \([a, b] \subset \mathbf{R}\) it is necessary and sufficient that it be convex on \(]a, b[\) and that one has \(f(a) \geqslant f(a+)\) and \(f(b) \geqslant f(b-)\) .
\item
  Let \(f\) be a convex function on an open interval \(]a, +\infty[\) ; if there exists a point \(c > a\) such that \(f\) is strictly increasing on \(]c, +\infty[\) then \(\lim_{x \to +\infty} f(x) = +\infty\) .
\item
  Let \(f\) be a convex function on an interval \(]a, +\infty[\) ; show that \(f(x)/x\) has a limit (finite or equal to \(+\infty\) ) as \(x\) tends to \(+\infty\) ; this limit is also that of \(f_d'(x)\) and of \(f_g'(x)\) ; it is \(>0\) if \(f(x)\) tends to \(+\infty\) as \(x\) tends to \(+\infty\) .
\item
  Let \(f\) be a convex function on the interval \(]a, b[\) where \(a \geqslant 0\) ; show that on this interval the function \(x \mapsto f(x) - xf'(x)\) (the ``ordinate at the origin'' of the right semi-tangent at the point \(x\) to the graph of \(f\) ) is decreasing (strictly decreasing if \(f\) is strictly convex).
\end{enumerate}

Deduce that:

\(a)\) If \(f\) admits a finite right limit at the point \(a\) then \((x - a)f_d'(x)\) has a right limit equal to 0 at this point.

\begin{enumerate}
\def\labelenumi{\alph{enumi})}
\setcounter{enumi}{1}
\item
  On {]}a, b{[} either \(f(x)/x\) is increasing, or \(f(x)/x\) is decreasing, or else there exists a \(c \in ]a\) , b{[} such that \(f(x)/x\) is decreasing on {]}a, c{[} and increasing on {]}c, b{[}.
\item
  Suppose that \(b = +\infty\) : show that if
\end{enumerate}

\[
\beta = \lim _ {x \rightarrow + \infty} \left(f (x) - x f _ {d} ^ {\prime} (x)\right)
\]

is finite, then so is \(\alpha = \lim_{x \to +\infty} f(x) / x\) , and that the line \(y = \alpha x + \beta\) is asymptotic \(^{5}\) to the graph of f, and lies below this graph (strictly below if f is strictly convex).

\begin{enumerate}
\def\labelenumi{\arabic{enumi})}
\setcounter{enumi}{7}
\tightlist
\item
  Let \(f\) be a finite real function, upper semi-continuous on an open interval \(I \subset \mathbf{R}\) . Then \(f\) is convex if and only if \(\limsup_{h \to 0, h \neq 0} \frac{f(x + h) + f(x - h) - 2f(x)}{h^2} \geqslant 0\) for all \(x \in I\) . (First show that, for all \(\varepsilon > 0\) the function \(f(x) + \varepsilon x^2\) is convex, using prop. 9 of I, p.~31.)
\end{enumerate}

¶ 9) Let \(f\) be a finite real function, lower semi-continuous on an interval \(\mathrm{I} \subset \mathbf{R}\) . For \(f\) to be convex on \(\mathrm{I}\) it suffices that, for every pair of points \(a, b\) of \(\mathrm{I}\) such that \(a < b\) there exists one point \(z\) such that \(a < z < b\) , and that \(\mathbf{M}_z\) be below the segment \(\mathbf{M}_a\mathbf{M}_b\) (argue by contradiction, noting that the set of points \(x\) such that \(\mathbf{M}_x\) lies strictly above \(\mathbf{M}_a\mathbf{M}_b\) is open).

¶ 10) Let \(f\) be a finite real function defined on an interval \(\mathrm{I}\subset \mathbf{R}\) , such that

\[
f \left(\frac {x + y}{2}\right) \leqslant \frac {1}{2} (f (x) + f (y))
\]

for all x, y in I. Show that if f is bounded above on one open interval {]}a, b{[} contained in I, then f is convex on I (show first that f is bounded above on every compact interval contained in I, then that f is continuous at every interior point of I).

¶ 11) Let \(f\) be a continuous function on an open interval \(\mathrm{I} \subset \mathbf{R}\) , having a finite right derivative at every point of \(\mathrm{I}\) . If for every \(x \in \mathrm{I}\) and every \(y \in \mathrm{I}\) such that \(y > x\) the point \(\mathbf{M}_y = (y, f(y))\) lies above the right semi-tangent to the graph of \(f\) at the point \(\mathbf{M}_x = (x, f(x))\) , show that \(f\) is convex on \(\mathrm{I}\) (using the mean value theorem show that \(f(y) = f(x)\) ).

\[
f _ {d} ^ {\prime} (y) \geqslant \frac {f (y) - f (x)}{y - x} \text {   for   } x <   y).
\]

Give an example of a function which is not convex, has a finite right derivative everywhere, and such that for every \(x \in I\) there exists a number \(h_{x} > 0\) depending on x such that \(M_{y}\) lies above the right semi-tangent at the point \(M_{x}\) for all y such that \(x \leqslant y \leqslant x + h_{x}\) . This last condition is nevertheless sufficient for f to be convex, if one supposes further that f is differentiable on I (use I, p.~12, corollary).

¶ 12) Let \(f\) be a continuous real function on an open interval \(\mathrm{I}\subset \mathbf{R}\) ; suppose that for every pair \((a,b)\) of points of I such that \(a < b\) the graph of \(f\) lies either entirely above or entirely below the segment \(\mathbf{M}_a\mathbf{M}_b\) on the interval \([a,b]\) . Show that \(f\) is convex on all of I or concave on all I (if in \(]a\) , \(b[\) there is a point \(c\) such that \(\mathbf{M}_c\) lies strictly above the segment \(\mathbf{M}_a\mathbf{M}_b\) show that for every \(x\in \mathrm{I}\) such that \(x > a\) the graph of \(f\) lies above the segment \(\mathbf{M}_a\mathbf{M}_x\) on the interval \([a,x]\) ).

\begin{enumerate}
\def\labelenumi{\arabic{enumi})}
\setcounter{enumi}{12}
\item
  Let \(f\) be a differentiable real function on an open interval \(\mathrm{I} \subset \mathbf{R}\) . Suppose that for every pair \((a, b)\) of points of \(\mathrm{I}\) such that \(a < b\) there exists a unique point \(c \in ]a, b[\) such that \(f(b) - f(a) = (b - a)f'(c)\) ; show that \(f\) is strictly convex on \(\mathrm{I}\) or strictly concave on \(\mathrm{I}\) (show that \(f'\) is strictly monotone on \(\mathrm{I}\) ).
\item
  Let \(f\) be a convex real function and strictly monotone on an open interval \(I \subset \mathbf{R}\) ; let \(g\) be the inverse function of \(f\) (defined on the interval \(f(I)\) ). Show that if \(f\) is decreasing (resp. increasing) on \(I\) , then \(g\) is convex (resp. concave) on \(f(I)\) .
\item
  Let I be an interval contained in \(]0, +\infty[\) ; show that if \(f(1 / x)\) is convex on I then so is \(xf(x)\) , and conversely.
\end{enumerate}

* 16) Let \(f\) be a positive convex function on \(]0, +\infty[\) , and \(a, b\) two arbitrary real numbers. Show that the function \(x^a f(x^{-b})\) is convex on \(]0, +\infty[\) in the following cases:

\[
1 ^ {\circ} a = \frac {1}{2} (b + 1), | b | \geqslant 1;
\]

\(2^{\circ} x^{a} f(x^{-b})\) is increasing, \(a(b - a) \geqslant 0\) , \(a \geqslant \frac{1}{2}(b + 1)\) ;

\(3^{\circ} x^{a} f(x^{-b})\) is decreasing, \(a(b-a) \geqslant 0\) , \(a \leqslant \frac{1}{2}(b+1)\) .

Under the same hypotheses on \(f\) show that \(e^{x / 2}f(e^{-x})\) is convex (use exerc. 2 of I, p.~45).*

\begin{enumerate}
\def\labelenumi{\arabic{enumi})}
\setcounter{enumi}{16}
\item
  Let \(f\) and \(g\) be two positive convex functions on an interval \(I = [a, b]\) ; suppose that there exists a number \(c \in \mathrm{I}\) such that in each of the intervals \([a, c]\) and \([c, b]\) the functions \(f\) and \(g\) vary in the same sense. Show that the product \(fg\) is convex on \(\mathrm{I}\) .
\item
  Let \(f\) be a convex function on an interval \(\mathrm{I} \subset \mathbf{R}\) and \(g\) a convex increasing function on an interval containing \(f(\mathrm{I})\) ; show that \(g \circ f\) is convex on \(\mathrm{I}\) .
\end{enumerate}

¶ 19) Let \(f\) and \(g\) be two finite real functions, \(f\) being defined and continuous on an interval I, and \(g\) defined and continuous on \(\mathbf{R}\) . Suppose that for every pair \((\lambda, \mu)\) of real numbers the function \(g(f(x) + \lambda x + \mu)\) is convex on I.

\begin{enumerate}
\def\labelenumi{\alph{enumi})}
\item
  Show that \(g\) is convex and monotone on \(\mathbf{R}\) .
\item
  If g is increasing (resp. decreasing) on R, show that f is convex (resp. concave) on I (use prop. 7).
\end{enumerate}

\begin{enumerate}
\def\labelenumi{\arabic{enumi})}
\setcounter{enumi}{19}
\item
  Show that the set \(\mathfrak{K}\) of convex functions on an interval \(\mathrm{I} \neq \mathbf{R}\) is reticulated for the order '' \(f(x) \leqslant g(x)\) for every \(x \in \mathrm{I}"\) (Set Theory, III, p.146). Give an example of two convex functions \(f\) , \(g\) on \(\mathrm{I}\) such that their infimum in \(\mathfrak{K}\) takes a value different from \(\inf(f(x), g(x))\) at certain points. Give an example of an infinite family \((f_{\alpha})\) of functions in \(\mathfrak{K}\) such that \(\inf_{\alpha} f_{\alpha}(x)\) is finite at every point \(x \in \mathrm{I}\) and yet there is no function in \(\mathfrak{K}\) less than all the \(f_{\alpha}\) .
\item
  Let \(f\) be a finite real function, upper semi-continuous on an open interval \(\mathbf{I} \subset \mathbf{R}\) . For \(f\) to be strictly convex on \(\mathbf{I}\) it is necessary and sufficient that there be no line locally above the graph \(\mathbf{G}\) of \(f\) at a point of \(\mathbf{G}\) .
\item
  Let \(f_{1}, \ldots, f_{n}\) be continuous convex functions on a compact interval \(I \subset \mathbf{R}\) ; suppose that for all \(x \in I\) one has \(\sup(f_{j}(x)) \geqslant 0\) . Show that there exist \(n\) numbers \(\alpha_{j} \geqslant 0\) such that \(\sum_{j=1}^{n} \alpha_{j} = 1\) and that \(\sum_{j=1}^{n} \alpha_{j} f_{j}(x) \geqslant 0\) on \(I\) . (First treat the case \(n = 2\) , considering a
\end{enumerate}

point \(x_{0}\) where the upper envelope \(\sup(f_{1}, f_{2})\) attains its minimum; when \(x_{0}\) is interior to I determine \(\alpha_{1}\) so that the left derivative of \(\alpha_{1}f_{1} + (1 - \alpha_{1})f_{2}\) is zero at \(x_{0}\) . Pass to the general case by induction on n; use the induction hypothesis for the restrictions of \(f_{1}, \ldots, f_{n-1}\) to the compact interval where \(f_{n}(x) \leqslant 0\) .

\begin{enumerate}
\def\labelenumi{\arabic{enumi})}
\setcounter{enumi}{22}
\item
  Let \(f\) be a continuous real function on a compact interval \(\mathrm{I} \subset \mathbf{R}\) ; among the functions \(g \leqslant f\) which are convex on \(\mathrm{I}\) there exists one, \(g_0\) , larger than all the others. Let \(\mathrm{F} \subset \mathrm{I}\) be the set of \(x \in \mathrm{I}\) where \(g_0(x) = f(x)\) ; show that \(\mathrm{F}\) is not empty and that on each of the open intervals contiguous to \(\mathrm{F}\) the function \(g_0\) is equal to an affine linear function (argue by contradiction).
\item
  Let \(\mathrm{P}(x)\) be a polynomial of degree \(n\) with real coefficients all of whose roots are real and contained in the interval \([-1, 1]\) . Let \(k\) be an integer such that \(1 \leqslant k \leqslant n\) . Show that the rational function
\end{enumerate}

\[
f (x) = x + \frac {\mathrm{P} ^ {(k - 1)} (x)}{\mathrm{P} ^ {(k)} (x)}
\]

is increasing on every interval of R on which it is defined; if \(c_{1} < c_{2} < \cdots < c_{r}\) are its poles (contained in \([-1, 1]\) ), then f is convex for \(x < c_{1}\) and concave for \(x > c_{r}\) . Deduce that when a runs through \([-1, 1]\) the length of the largest interval containing the zeros of the \(k^{th}\) derivative of \((x - a)\mathrm{P}(x)\) attains its largest value when a = 1 or a = -1.

\begin{enumerate}
\def\labelenumi{\arabic{enumi})}
\setcounter{enumi}{24}
\tightlist
\item
  One says that a real function \(f\) defined on \([0, +\infty[\) is superadditive if one has \(f(x + y) \geqslant f(x) + f(y)\) for \(x \geqslant 0, y \geqslant 0\) , and if \(f(0) = 0\) .
\end{enumerate}

\begin{enumerate}
\def\labelenumi{\alph{enumi})}
\item
  Give examples of discontinuous superadditive functions.
\item
  Show that every convex function \(f\) on \([0, +\infty[\) such that \(f(0) = 0\) is superadditive.
\item
  If \(f_{1}\) and \(f_{2}\) are superadditive then so is \(\inf (f_1, f_2)\) ; using this, exhibit examples of nonconvex continuous superadditive functions.
\item
  If f is continuous and \(\geqslant0\) on an interval \([0,a]\) (a\textgreater0), such that \(f(0)=0\) and \(f(x/n)\leqslant f(x)/n\) for each integer \(n\geqslant1\) , show that f has a right derivative at the point 0 (argue by contradiction). In particular, every continuous superadditive function which is \(\geqslant0\) admits a right derivative at the point 0.
\end{enumerate}

